\section{Cyclotomic-adapted inertia generators on the actual hosts}
\label{sec:cubic-frame-generators}

The local generator clause in the source is stronger than a count of
the rank of local units: all but one generator must lie in the kernel
of the full cyclotomic character.  The following verifies that clause
on the actual selected-place hosts, simultaneously for every finite
residual character under consideration.  Its global exponent assertion
is a qualitative class-field consequence, not a class-rank estimate.

\begin{proposition}[AC-G99-frame: the actual local generators and global exponent]
\label{AC-G99-frame}
Let \(p=269\), let \(F\) be an actual selected-place host at a finite
cyclotomic layer as in AC-G93-fill, and let \(H=KF\) be its CM host.
For every \(v\mid p\), put \(E=F_v\), \(d_v=[E:\mathbb Q_p]\),
and retain the full absolute group \(G_E\), its inertia \(I_E\),
and its full cyclotomic character \(\epsilon:G_E\to\mathbb Z_p^\times\).
Let \(\mathcal T_v\) be any finite family of continuous characters
of \(G_E\) with finite image of order prime to \(p\).  It can include
all actual residual tuple characters and the finite residual rational character
\(\gamma\).

There exist \(u_{v,0},\ldots,u_{v,d_v-1}\in I_E\) such that
\[
 \tau(u_{v,i})=1\quad(\tau\in\mathcal T_v),\qquad
 \epsilon(u_{v,i})=1\quad(1\le i<d_v),                 \tag{G99.1}
\]
and, for every finite abelian \(p\)-group \(A\) and continuous
homomorphism \(\nu:G_E\to A\),
\[
 \bigl(\nu(u_{v,i})=1\text{ for all }i\bigr)
       \ \Longrightarrow\ \nu(I_E)=1.                \tag{G99.2}
\]
Thus, with \(a_0=0\), these elements verify the literal source
\path{Frame.Ok.gen} whenever its residual character is one of
the specified characters.  The same elements work for all of them;
this assertion does not identify the two-root tuple with powers of
one character.

Let \(G=G_H/N\) be the actual joint residual-away-kernel quotient.
Choose one full local map at each member of every pair of places
above \(p\), as in the source's \path{loc2}.  There is a positive
integer \(h_{\rm cl}\), depending on this fixed host, such that every
continuous homomorphism \(\nu:G\to A\) to a finite abelian
\(p\)-group which kills their inertia satisfies
\[
                         \nu(g)^{h_{\rm cl}}=1
                         \quad(g\in G).              \tag{G99.3}
\]
One may take the exponent of
\(\operatorname{Cl}(\mathcal O_H)_{(p)}\), with exponent one for
the zero group.  This verifies the literal \path{Frame.Ok.cl}
field.  Neither assertion bounds \(h_{\rm cl}\) or any cubic
strict class-sector dimension uniformly over varying hosts.
\end{proposition}

\begin{proof}
\emph{1. The actual local unit lattice.}
The coefficient-prime completions arise from unramified extensions
of \(\mathbb Q_p\) of degree one or two by the unramified
\(p\)-power local extensions allowed in AC-G83.  The G92
selected-place extensions split completely there, and the cyclotomic
layers have local \(p\)-power degree.  In particular \([E:\mathbb Q_p]\) has
no factor \(67\).  If \(E\) contained \(\mu_p\), the local degree
\([\mathbb Q_p(\mu_p):\mathbb Q_p]=268=4\cdot67\) would divide
\([E:\mathbb Q_p]\).  Thus \(\mu_p(E)=1\).

Write \(U^1=1+\mathfrak m_E\).  It is an abelian pro-\(p\)
group: its successive unit-filtration quotients are additive
groups of the finite residue field.  For an integer
\(n>e(E/\mathbb Q_p)/(p-1)\), logarithm and exponential identify
\[
                       U^n\simeq\mathfrak m_E^n.     \tag{G99.4}
\]
For completeness, convergence follows by comparing the valuations
of \(x^m/m\) and \(x^m/m!\); the strict displayed inequality makes
their valuations tend to infinity.  The inverse formal identities
therefore hold on these neighborhoods.  The additive lattice
\(\mathfrak m_E^n\) is free of rank \(d_v\) over \(\mathbb Z_p\).
The finite-index extension \(U^1/U^n\) is a finite \(p\)-group,
so \(U^1\), with its continuous power action, is a finitely generated
\(\mathbb Z_p\)-module of rank \(d_v\).  It has no nonzero torsion:
a torsion element would be a \(p\)-power root of unity and would
give a primitive \(p\)th root after taking a power.  The structure
theorem over \(\mathbb Z_p\) gives
\[
                            U^1\simeq\mathbb Z_p^{d_v}. \tag{G99.5}
\]

Local reciprocity identifies \(G_E^{\rm ab}\) with the profinite
completion of \(E^\times\), and the image of inertia with
\(\mathcal O_E^\times\).  In particular every element of \(U^1\)
has an inertia lift.  These are the local reciprocity and existence
theorems and their unit/inertia consequence, retained as inputs;
they are not newly proved here.

\emph{2. Adapt a basis to the full cyclotomic character.}
On \(U^1\), cyclotomy gives a continuous \(\mathbb Z_p\)-linear
homomorphism, in multiplicative notation,
\[
                  e_{\rm cyc}:U^1\longrightarrow1+p\mathbb Z_p.
                                                               \tag{G99.6}
\]
Its image is infinite and open.  Indeed \(I_E\) is a finite-index
subgroup of \(I_{\mathbb Q_p}\) inside the fixed local algebraic
closure, and the image of \(I_{\mathbb Q_p}\) under cyclotomy is
\(\mathbb Z_p^\times\), since the cyclotomic extensions of
\(\mathbb Q_p\) are totally ramified.  Its image on \(I_E\)
therefore has finite index.  Passing from
\(\mathcal O_E^\times\) to its finite-index subgroup \(U^1\)
still gives finite index in that image.  A pro-\(p\) group maps
into \(1+p\mathbb Z_p\), proving the assertion.

The image in (G99.6) is a free rank-one \(\mathbb Z_p\)-module.
The exact sequence
\[
 0\longrightarrow\ker e_{\rm cyc}\longrightarrow U^1
   \longrightarrow\operatorname{im}e_{\rm cyc}\longrightarrow0
                                                               \tag{G99.7}
\]
splits over \(\mathbb Z_p\).  Choose one basis vector lifting
a generator of the image and a basis of its kernel, which is free
of rank \(d_v-1\).  Lift all these basis vectors through inertia.
These lifts are the \(u_{v,i}\), and the kernel basis gives the
second assertion in (G99.1) exactly, for the full cyclotomic
character rather than its reduction.

Each character in \(\mathcal T_v\) factors through
\(G_E^{\rm ab}\).  Its restriction to the pro-\(p\) group \(U^1\)
is trivial because its finite image has order prime to \(p\).
This gives the first assertion in (G99.1).  For \(\nu\) in
(G99.2), its inertia restriction factors through
\(\mathcal O_E^\times\).  The Teichmuller residue-unit factor
has order prime to \(p\) and is killed by \(\nu\).
The basis vectors topologically generate \(U^1\); continuity
therefore makes vanishing on them equivalent to vanishing on
all of \(U^1\).  This proves (G99.2).  The cyclotomic kernel
condition places every \(u_{v,i}\), \(i\ne0\), in the literal
source subgroup \(N_v=I_v\cap\ker\epsilon\).

\emph{3. The global exponent.}
Inflate \(\nu\) to \(G_H\).  At an away-\(p\) place, its inertia
image in \(G\) factors through the finite image of the residual
tuple.  That image has order prime to \(p\), so its homomorphism
to \(A\) is trivial.  The assumed conditions at both members of
each coefficient-prime pair kill all inertia above \(p\).
Thus the inflated character is unramified at every finite place.
The field \(H\) is CM, so there are no real places to modify.
Global reciprocity identifies the maximal everywhere-unramified
abelian pro-\(p\) quotient with
\(\operatorname{Cl}(\mathcal O_H)_{(p)}\).
Its finite exponent gives (G99.3), uniformly for every target \(A\)
and character \(\nu\) on this fixed host.
\end{proof}

\paragraph{Exact source and geometric scope.}
The literal \path{gen}, \path{cl}, and \path{Nv} declarations are
in the complete inspected
\path{FrobeniusSmoothVectors/ContinuousCoboundaries2.lean},
at pinned commit
\texttt{7cbf7a0c32fca1badf05931879a37c3f8a2b03b4}.
For local reciprocity, the exact retained reference is Milne,
\emph{Class Field Theory}, version 3.1 (6 May 1997), I,
Theorems 1.1--1.2, printed p.\ 10, and their unit/inertia
and completion consequences on printed pp.\ 11--13.
For the global assertion use V, Theorems 3.5--3.6 and
Example 3.9, printed pp.\ 123--124.  The applied statements
and local completion argument were reread in this run.
The logarithm argument and basis adaptation above prove the
specific generator form needed here; a new proof of reciprocity
is not claimed.

The inertia generators describe the formal character torus at the
coefficient-prime points of the arithmetic curve, with one
cyclotomic direction and \(d_v-1\) kernel directions.
The exponent is the finite-etale unramified-cover statement in
the already checked SGA~1 dictionary.  Neither is a new
Grothendieck archive attribution.  The actual two-root dimension
and propagation theorems must still be checked individually;
this proposition does not supply a numerical adjacent-sector
bound, a prescribed varying-character Hecke curve, or H04.4.
